% Generated draft: requires mathematical and source review.
\documentclass[11pt]{article}
\usepackage[margin=25mm]{geometry}
\usepackage{fontspec}
\setmainfont{Latin Modern Roman}
\usepackage{amsmath,amssymb,mathtools,unicode-math}
\setmathfont{Latin Modern Math}
\usepackage{enumitem,needspace,xurl,hyperref}
\setlength{\parindent}{0pt}
\setlength{\parskip}{5pt}
\setlength{\emergencystretch}{3em}
\newcommand{\N}{\mathbb{N}}
\newcommand{\Z}{\mathbb{Z}}
\newcommand{\Q}{\mathbb{Q}}
\newcommand{\R}{\mathbb{R}}
\newcommand{\C}{\mathbb{C}}
\newcommand{\F}{\mathbb{F}}
\newcommand{\A}{\mathbb{A}}
\newcommand{\PP}{\mathbb{P}}
\newcommand{\E}{\mathbb{E}}
\newcommand{\eps}{\varepsilon}
\newcommand{\dd}{\,\mathrm{d}}
\begin{document}
\section*{UnsolvedMath ID: 30003439}
\textbf{Source problem code:} catalogue entry
\subsection{Definitions and mathematical statement}
The supplied record describes an \emph{algebraic convex cone} as a convex cone defined by homogeneous multivariate polynomial inequalities. It does not specify the coefficient field, the allowed inequality signs, or whether the cone includes its boundary. These conventions must be fixed before the question has a fully precise formulation.

For reference, under the usual finite-dimensional real convention, a nonzero homogeneous polynomial \(p\in\R[x_1,\ldots,x_n]\) is \emph{hyperbolic with respect to} \(e\in\R^n\) if \(p(e)>0\) and, for every \(x\in\R^n\), every root of \(t\mapsto p(x-te)\) is real. Its open hyperbolicity cone is the connected component containing \(e\) of \(\{x\in\R^n:p(x)>0\}\); its closed hyperbolicity cone is the closure of that component. The supplied record does not specify which boundary convention is intended.

\textbf{Existence question.} Does there exist an algebraic convex cone \(K\), defined by homogeneous multivariate polynomial inequalities, which admits a normal barrier of negative curvature, but which is not a hyperbolicity cone? Here the last condition means that no homogeneous hyperbolic polynomial, with any hyperbolicity direction, has \(K\) as its hyperbolicity cone, using the intended boundary convention consistently.

\textbf{Unresolved analytic qualifications.} The supplied record gives no definition of a \emph{normal barrier} or of \emph{negative curvature}. It therefore does not determine the barrier's domain, regularity, normalization, parameters, or required derivative inequalities. It also does not state whether \(K\) must be closed, pointed, or full-dimensional, or whether barriers are taken on an interior or a relative interior. No additional hypotheses or analytic definition are imposed in this draft. The existence question above preserves the supplied target, but a complete quantified analytic formulation requires these missing conventions.
\subsection{Short English statement}
The question asks whether a convex cone described by homogeneous polynomial inequalities can admit a normal barrier of negative curvature while admitting no representation as a hyperbolicity cone. The meanings of the barrier conditions and the cone boundary conventions require clarification. The supplied record labels the problem open; this draft does not independently verify that status.
\subsection{Sources}
\begin{itemize}
\item Local UnsolvedMath registry, supplied record with permanent ID \(30003439\), title \emph{Negatively Curved Barriers Beyond Hyperbolicity Cones}; collection: \emph{Oberwolfach Reports Open Problems}; exact source problem code: \texttt{catalogue entry}. The supplied proposed year is 2017; no individual proposer is identified.
\item Supplied report citation: \emph{Real Algebraic Geometry With a View Toward Moment Problems and Optimization} (2018), Oberwolfach Reports, pp.~771--862, DOI \texttt{10.4171/owr/2017/14}. \url{https://doi.org/10.4171/owr/2017/14}. The background describes Oberwolfach Reports as documenting mathematical workshops and their open-problem discussions. No exact problem page or quoted definitions are supplied.
\item Supplied literature source: Netzer--Sanyal, smooth hyperbolicity cones. \url{https://arxiv.org/abs/1208.0441}.
\item Supplied literature source: Oberwolfach hyperbolicity-cone report. \url{https://doi.org/10.4171/owr/2017/14}.
\item Literature assessment dated exactly \texttt{2026-08-21}: the supplied curation labels the problem open and states that hyperbolicity cones supply logarithmic barriers with negative curvature, that related Helton--Nie statements have counterexamples, and that the narrower existence question for an algebraic non-hyperbolicity cone with such a barrier remains unresolved. This is supplied dated curation derived from the OWR material, not an independently verified literature assessment; its time-sensitive status and cited claims require independent verification.
\end{itemize}
\end{document}
